% TOP_PROBLEM: {"id": 20000949, "title": "Loose hypergraph triangles versus cliques: an edge-rooted component certificate", "category_id": 2, "category": {"id": 2, "name": "combinatorics", "display_name": "Combinatorics", "description": "Counting problems, graph theory, discrete structures.", "slug": "combinatorics", "order_index": 2, "created_at": "2026-07-31T15:26:25.671Z"}, "status": "partially_solved", "problem_number": "AIM-COMBINATORICS-0074", "set_id": 14}
% FIRST_POSTED: 2026-10-01
% ENTRY_KIND: statement_only
% UnsolvedMath ID 20000949; source catalogue code AIM-COMBINATORICS-0074
% !TeX program = lualatex
% Definition and sources only; this file is not an LLM solution attempt.
\documentclass[11pt,a4paper]{article}
\usepackage[margin=25mm]{geometry}
\usepackage{fontspec}
\setmainfont{lmroman10-regular.otf}[BoldFont=lmroman10-bold.otf,
 ItalicFont=lmroman10-italic.otf,BoldItalicFont=lmroman10-bolditalic.otf]
\usepackage{amsmath,amssymb,mathtools}
\usepackage{unicode-math}
\setmathfont{latinmodern-math.otf}
\AtBeginDocument{\let\setminus\smallsetminus}
\usepackage{xurl,hyperref}
\hypersetup{colorlinks=true,urlcolor=blue}
\setcounter{secnumdepth}{0}
\setlength{\parindent}{0pt}
\setlength{\parskip}{5pt}
\setlength{\emergencystretch}{3em}
\begin{document}
\section*{Loose hypergraph triangles versus cliques: an edge-rooted component certificate}\label{20000949}
{\small UnsolvedMath ID 20000949 · AIM-COMBINATORICS-0074 · Combinatorics · AIM Workshop Problem Lists}
\subsection{Definitions and mathematical statement}
Let $T$ be the $3$-uniform hypergraph with vertex set $[6]$ and edge
set
\[
              \{\{1,2,3\},\{3,4,5\},\{5,6,1\}\}.
\]
It is a loose triangle: each pair of edges meets in one vertex, and
the three intersection vertices are distinct. In particular, it is
not the tight three-edge path or a triangle in an ordinary graph.
For $t\ge3$, define $R_T(t)=r_3(T,K_t^{(3)})$ as the least $N$ for
which every red--blue coloring of $\binom{[N]}3$ contains a red copy
of $T$ or a blue complete $3$-uniform hypergraph on $t$ vertices.
The red copy uses six distinct vertices and need not be induced.

The proposed improvement is
\[
                         R_T(t)=o(t^{3/2})\qquad(t\to\infty).
\]
Precisely, for every real $\varepsilon>0$ there should be a threshold
$t_0$ such that $R_T(t)\le\varepsilon t^{3/2}$ for all $t\ge t_0$.
Equivalently, every sufficiently large $T$-free triple system on
$N$ vertices should have an independent set larger than
$C N^{2/3}$, for each fixed $C>0$ once $N$ is sufficiently large.
An independent set is a set containing no hyperedge. This reformulation
uses the red triples as the hypergraph and the blue clique as an
independent set. The coefficient is required to improve without bound;
a fixed constant improvement to $O(t^{3/2})$ would not suffice.
\subsection{Short English statement}
Is the Ramsey number of a loose triple-system triangle against a clique little-oh of the clique size to the power three halves?
\subsection{Sources}
\begin{itemize}
\item[S1] ULAM.AI and the UnsolvedMath Contributors, supplied problems.json, record 20000949 (AIM-COMBINATORICS-0074), AIM Workshop Problem Lists. Catalogue identity and source transcription; CC BY 4.0.
\url{https://huggingface.co/datasets/ulamai/UnsolvedMath}.
\item[S2] American Institute of Mathematics, Graph Ramsey theory, item 27. Original workshop question underlying AIM-COMBINATORICS-0074.
\url{https://www.overleaf.com/read/mnvcscjjysvg}.
\end{itemize}
\end{document}
